\documentclass[10pt]{article}
\usepackage[a4paper,margin=0.45in]{geometry}
\usepackage{amsmath,amssymb,amsfonts,mathtools}
\usepackage[dvipsnames]{xcolor}
\usepackage{booktabs,array}
\usepackage[colorlinks=true,linkcolor=blue]{hyperref}
\usepackage{etoolbox}
\AtBeginEnvironment{align}{\small}
\allowdisplaybreaks
\numberwithin{equation}{section}
\setlength{\emergencystretch}{3em}
\newcommand{\new}[1]{{\color{blue}#1}}
\newcommand{\old}[1]{{\color{red}#1}}
\newcommand{\LO}{{\rm LO}}
\newcommand{\MSb}{\overline{\rm MS}}
\title{Group II Row IV: the compact form, the running coupling,\\ and why line 2 is not a running-coupling term}
\date{}
\begin{document}
\maketitle
\noindent Conventions: $a\equiv\frac{k^+}{\vee}$ (kept exact), $\alpha_s=\frac{g^2}{4\pi}$, so $\frac{g^2}{8\pi^2}=\frac{\alpha_s}{2\pi}$ and $\frac{g^2N_c}{4\pi^2}=\frac{\alpha_s}{2\pi}2N_c$. $\beta_0=\frac{11}{3}N_c$ (no quark loops in this row). $c_0=2e^{-\gamma_E}$. $d_\perp=2-\epsilon$. $\mathbf k\neq0$ (measured gluon).
\new{Blue} = added/changed, \old{red} = removed, relative to your compact form.

\section{Definitions}
Your $\frac{d^3N_{\xi,\rm LO}}{d^3k}$ with a weight $f$ placed \emph{inside} the integral:
\begin{align}
\LO_\xi[f]&\equiv\frac{1}{(2\pi)^{3}}\frac{g^{2}}{2\pi^{2}}\frac{1}{k^{+}}\int_{x,y,x',y'}e^{-ik\cdot(y'-y)\frac1\xi}\,f\,\frac{(x-y)\cdot(x'-y')}{(x-y)^{2}(x'-y')^{2}}\left[U^{e'c}(y')-U^{e'c}(x')\right]\left[U^{ec}(y)-U^{ce}(x)\right]\rho^{e'}(x')\rho^{e}(x),\label{eq:LOxi}\\
\LO^{(y)}_\xi[f]&\equiv\frac{1}{(2\pi)^{3}}\frac{g^{2}}{2\pi^{2}}\frac{1}{k^{+}}\int_{x,y,x',y'}e^{-ik\cdot(y'-y)\frac1\xi}\,f\,\frac{(x-y)\cdot(x'-y)}{(x-y)^{2}(x'-y)^{2}}\left[U^{e'c}(y)-U^{e'c}(x')\right]\left[U^{ec}(y)-U^{ce}(x)\right]\rho^{e'}(x')\rho^{e}(x),\label{eq:LOy}\\
\LO^{\rm lp}_\xi&\equiv\frac{1}{(2\pi)^{3}}\frac{g^{2}}{2\pi^{2}}\frac{1}{k^{+}}\int_{x,y,x',y'}e^{-ik\cdot(y'-y)\frac1\xi}\Bigg[\frac{(x-y)\cdot(x'-y')}{(x-y)^{2}(x'-y')^{2}}-d_\perp\frac{(x-y)\cdot(y'-y)\,(x'-y')\cdot(y'-y)}{(x-y)^{2}(x'-y')^{2}(y'-y)^{2}}\Bigg]\notag\\
&\qquad\times\left[U^{e'c}(y')-U^{e'c}(x')\right]\left[U^{ec}(y)-U^{ce}(x)\right]\rho^{e'}(x')\rho^{e}(x).\label{eq:LOlp}
\end{align}
$\LO_\xi\equiv\LO_\xi[1]=\frac1\xi\frac{d^3N_{\rm LO}}{d^3k}\big(\frac{k^+}{\xi},\frac{\mathbf k}{\xi}\big)$ and $\LO\equiv\LO_{\xi=1}=\frac{d^3N_{\rm LO}}{d^3k}$. $\LO^{(y)}_\xi$ is $\LO_\xi$ at $y'=y$ in the kernel and in $U(y')$ (your line 2). $H_4$ is your line 3 without $-(1-\frac{k^+}{\vee})$:
\begin{align}
H_4\equiv\frac{1}{(2\pi)^{3}}\frac{g^{4}N_{c}}{8\pi^{5}}\frac{1}{k^{+}}\int_{x,x',z,w,w'}e^{-ik\cdot(w'-w)}\frac{(w-z)\cdot(w'-z)}{(w-z)^{2}(w'-z)^{2}}\frac{(x-z)\cdot(x'-z)}{(x-z)^{2}(x'-z)^{2}}\left[U^{e'c}(z)-U^{e'c}(x')\right]\left[U^{ec}(z)-U^{ce}(x)\right]\rho^{e'}(x')\rho^{e}(x).\label{eq:H4}
\end{align}
\begin{align}
\ell_\Delta(\mu)&\equiv2\gamma_E+\log\frac{(y'-y)^2\mu^2}{4}=\log\frac{(y'-y)^2\mu^2}{c_0^2},\label{eq:ell}\\
P_{gg}(\xi)&=2N_c\Big[\frac{\xi}{(1-\xi)_+}+\frac{1-\xi}{\xi}+\xi(1-\xi)\Big]+\frac{\beta_0}{2}\delta(1-\xi),\qquad
P^{\rm real}_{gg}(\xi)\equiv P_{gg}(\xi)-\frac{\beta_0}{2}\delta(1-\xi),\label{eq:P}
\end{align}
with the plus distribution of $[0,1]$: $\int_a^1d\xi\frac{g(\xi)}{(1-\xi)_+}=\int_a^1d\xi\frac{g(\xi)-g(1)}{1-\xi}+g(1)\log(1-a)$.

\section{Is the compact form correct?}
\begin{center}\small
\begin{tabular}{@{}l l p{0.62\textwidth}@{}}
\toprule
Line & Status & Comment\\\midrule
1 & \textbf{fix} & (i) The bracket depends on $y'-y$: it must be the weight $f$ \emph{inside} $\LO_\xi[f]$, not a factor outside $\frac{d^3N_{\xi,\rm LO}}{d^3k}$.\\
  & & (ii) The colour factor written after line 1 is already inside $\frac{d^3N_{\xi,\rm LO}}{d^3k}$ (your own definition): remove it.\\
  & & (iii) $\mathcal P_{gg}$ must be the full $P_{gg}$ of \eqref{eq:P}: standard plus distribution on $[0,1]$ (cut at $\xi\ge a$) and $+\frac{\beta_0}{2}\delta(1-\xi)$. Then $\mathcal P_{gg}-\frac{\beta_0}{2}\delta=P^{\rm real}_{gg}$ reproduces your original $\xi$-bracket \emph{and} your $\log(1-\frac{k^+}{\vee})$ line. With your prescription $[\frac1{p^+}]_+$ instead, add $-\log(1-a)\,\delta(1-\xi)$.\\
  & & (iv) $-2\gamma_E-\log(\dots)$ is the $\epsilon\to0$ form of $-\frac{d_\perp}{2}\ell_\Delta$; see \eqref{eq:B}.\\
2 & correct & Same as the original line. Its $-\log$ is $-\frac{d_\perp}{2}\log$ at $\epsilon\to0$ (no pole multiplies it). $2\gamma_E$ may be added or not: it multiplies $\delta^2(\mathbf k)$, see Sec.~\ref{sec:line2}.\\
3 & correct & $-(1-\frac{k^+}{\vee})H_4$.\\
4 & correct & $\frac{\alpha_s}{2\pi}2N_c\int_a^1d\xi\frac{1-\xi}{\xi}\LO^{\rm lp}_\xi$; $d_\perp\to2$ at the end (finite).\\
\bottomrule
\end{tabular}
\end{center}
The exact bracket, with $d_\perp=2-\epsilon$ (checked in \texttt{checks.py}):
\begin{equation}
-\frac{d_\perp}{\epsilon}-d_\perp\gamma_E-\frac{d_\perp}{2}\log\frac{(y'-y)^2\mu^2}{4}-1=-\frac2\epsilon-\frac{d_\perp}{2}\ell_\Delta(\mu)=-\frac2\epsilon-\ell_\Delta(\mu)+\frac\epsilon2\ell_\Delta(\mu).\label{eq:B}
\end{equation}
The last term has no pole in front of it and vanishes at $\epsilon\to0$.

\paragraph{Corrected compact form.}
\begin{align}
\frac{d^{3}N}{d^{3}k}={}&\frac{\alpha_s}{2\pi}\int_{a}^{1}\frac{d\xi}{\xi^{2}}\Big[\new{P_{gg}}(\xi)-\frac{\beta_{0}}{2}\delta(1-\xi)\Big]\,\new{\LO_\xi\Big[}-\frac{2}{\epsilon}-\new{\frac{d_\perp}{2}}\ell_\Delta(\mu_{\MSb})\new{\Big]}\;\old{\times\big[U^{e'c}(y')-U^{e'c}(x')\big]\big[U^{ec}(y)-U^{ce}(x)\big]\rho^{e'}(x')\rho^{e}(x)}\notag\\
&+\frac{\alpha_s}{2\pi}\,2N_c\int_a^1\frac{d\xi}{\xi^2}\frac{1-\xi}{\xi}\,\LO^{(y)}_\xi\Big[\frac{d_\perp}{2}\log\frac{(y'-y)^2\mu^2_{\MSb}}{4}\Big]
-\Big(1-\frac{k^+}{\vee}\Big)H_4+\frac{\alpha_s}{2\pi}2N_c\int_a^1d\xi\frac{1-\xi}{\xi}\LO^{\rm lp}_\xi.\label{eq:compact}
\end{align}

\section{Simplified: DGLAP $+$ running coupling $+$ finite}
\paragraph{Step 1: lines 1 and 2 as one $\xi$ integral.} Line 2 has the same $\frac{1-\xi}{\xi}$ as the small-$\xi$ part of $P_{gg}$, so
\begin{align}
\text{(1)}+\text{(2)}=\frac{\alpha_s}{2\pi}\int_a^1\frac{d\xi}{\xi^2}\Big\{&P^{\rm real}_{gg}(\xi)\,\LO_\xi\Big[-\frac2\epsilon-\frac{d_\perp}2\ell_\Delta(\mu)\Big]+2N_c\frac{1-\xi}{\xi}\,\LO^{(y)}_\xi\Big[\frac{d_\perp}{2}\ell_\Delta(\mu)\Big]\Big\}.\label{eq:s1}
\end{align}
For $\xi\to0$ the phase $e^{-ik\cdot(y'-y)/\xi}$ forces $y'\to y$, so the $\log$-part of $\frac{1-\xi}{\xi}\big(\LO_\xi-\LO^{(y)}_\xi\big)$ vanishes there: no $\log\frac{\vee}{k^+}$ is produced by \eqref{eq:s1} (it is already in BFKL).

\paragraph{Step 2: split $P^{\rm real}_{gg}=P_{gg}-\frac{\beta_0}2\delta(1-\xi)$.} At $\xi=1$, $\LO_\xi[f]=\LO[f]$:
\begin{align}
\eqref{eq:s1}=\underbrace{\frac{\alpha_s}{2\pi}\int_a^1\frac{d\xi}{\xi^2}\Big\{P_{gg}\,\LO_\xi\Big[-\frac2\epsilon-\frac{d_\perp}{2}\ell_\Delta(\mu)\Big]+2N_c\frac{1-\xi}{\xi}\LO^{(y)}_\xi\Big[\frac{d_\perp}{2}\ell_\Delta(\mu)\Big]\Big\}}_{\rm DGLAP}
+\underbrace{\frac{\alpha_s\beta_0}{4\pi}\,\LO\Big[\frac2\epsilon+\frac{d_\perp}{2}\ell_\Delta(\mu)\Big]}_{\rm running\ coupling}.\label{eq:s2}
\end{align}

\paragraph{Step 3: DGLAP.} The pole $\frac{\alpha_s}{2\pi}(-\frac2\epsilon)P_{gg}\otimes\LO$ is removed by the $\MSb$ fragmentation function, $D^{\rm bare}=D(\mu_F)+\frac{\alpha_s}{2\pi}\frac2\epsilon P_{gg}\otimes D$, with $\frac{d^3N_h}{d^3k}=\int\frac{d\zeta}{\zeta^3}D(\zeta)\frac{d^3N_g}{d^3k}(\frac{k}{\zeta})$. Setting $\mu=\mu_F$:
\begin{equation}
\boxed{\frac{d^3N}{d^3k}\bigg|_{\rm DGLAP}=-\frac{\alpha_s}{2\pi}\int_a^1\frac{d\xi}{\xi^2}\Big\{P_{gg}(\xi)\,\LO_\xi\Big[\ell_\Delta(\mu_F)\Big]-2N_c\frac{1-\xi}{\xi}\,\LO^{(y)}_\xi\Big[\ell_\Delta(\mu_F)\Big]\Big\}}\label{eq:dglap}
\end{equation}
($\frac{d_\perp}{2}\to1$: no pole is left in front). Check: $\mu_F^2\partial_{\mu_F^2}$ of \eqref{eq:dglap} is $-\frac{\alpha_s}{2\pi}P_{gg}\otimes\LO$, which cancels $\mu_F^2\partial_{\mu_F^2}D=\frac{\alpha_s}{2\pi}P_{gg}\otimes D$. The $\LO^{(y)}$ term contributes nothing to this derivative (Sec.~\ref{sec:line2}).

\paragraph{Step 4: running coupling} (details in Sec.~\ref{sec:rc}):
\begin{equation}
\boxed{\frac{d^3N}{d^3k}\bigg|_{\rm RC}=\frac{\alpha_s\beta_0}{4\pi}\,\LO\Big[\log\frac{(y'-y)^2\mu_R^2}{c_0^2}\Big]}\label{eq:rcbox}
\end{equation}

\paragraph{Step 5: finite.}
\begin{equation}
\boxed{\frac{d^3N}{d^3k}\bigg|_{\rm finite}=-\Big(1-\frac{k^+}{\vee}\Big)H_4+\frac{\alpha_s}{2\pi}\,2N_c\int_a^1d\xi\,\frac{1-\xi}{\xi}\,\LO^{\rm lp}_\xi\Big|_{d_\perp=2}}\label{eq:fin}
\end{equation}
So $\displaystyle\frac{d^3N_{\rm G2\,IV}}{d^3k}=\eqref{eq:dglap}+\eqref{eq:rcbox}+\eqref{eq:fin}$.

\section{Running coupling, step by step}\label{sec:rc}
\begin{enumerate}
\item \textbf{Source.} Only the $-\frac{\beta_0}{2}\delta(1-\xi)$ in line 1:
$\ \frac{\alpha_s}{2\pi}\big(-\frac{\beta_0}{2}\big)\LO\big[-\frac2\epsilon-\frac{d_\perp}{2}\ell_\Delta(\mu)\big]=\frac{\alpha_s\beta_0}{4\pi}\LO\big[\frac2\epsilon+\frac{d_\perp}2\ell_\Delta(\mu)\big]$.
\item \textbf{Coupling counterterm.} $\LO\propto g^2$. In $\MSb$, $\alpha_s^{\rm bare}\mu^{\epsilon}=\alpha_s(\mu_R)\big[1-\frac{\alpha_s\beta_0}{4\pi}\frac2\epsilon\big]$ (here $\frac1{\epsilon_{d=4-2\epsilon'}}=\frac2\epsilon$), so the LO term gives $-\frac{\alpha_s\beta_0}{4\pi}\frac2\epsilon\LO$.
\item \textbf{Sum} at $\mu=\mu_R$: the pole cancels, and $\frac{d_\perp}2\to1$:
$\ \frac{\alpha_s\beta_0}{4\pi}\LO\big[\ell_\Delta(\mu_R)\big]=\frac{\alpha_s\beta_0}{4\pi}\LO\big[\log\frac{(y'-y)^2\mu_R^2}{c_0^2}\big]$.
\item \textbf{Resummed.} With $\alpha_s(Q)=\alpha_s(\mu_R)\big[1-\frac{\alpha_s\beta_0}{4\pi}\log\frac{Q^2}{\mu_R^2}\big]$,
\begin{equation}
\LO\Big|_{\alpha_s(\mu_R)}+\frac{\alpha_s\beta_0}{4\pi}\LO\Big[\log\frac{(y'-y)^2\mu_R^2}{c_0^2}\Big]=\LO\Big|_{\alpha_s\to\alpha_s\left(\frac{c_0}{|y'-y|}\right)}+O(\alpha_s^3),\label{eq:rcres}
\end{equation}
i.e.\ the coupling of the emitted gluon runs with the amplitude--conjugate separation $|y'-y|$, inside the $y,y'$ integrals.
\item \textbf{Check.} $\mu_R^2\partial_{\mu_R^2}$: LO gives $-\frac{\alpha_s^2\beta_0}{4\pi}$, \eqref{eq:rcbox} gives $+\frac{\alpha_s^2\beta_0}{4\pi}$: $\mu_R$-independent.
\item With quarks, $\beta_0\to\frac{11}3N_c-\frac23n_f$; the $n_f$ part comes from quark-loop rows, not from this row.
\end{enumerate}
\noindent (This uses $\mu=\mu_{\MSb}$ as in your equation, i.e.\ $\gamma_E$ and $\log4\pi$ already in $\mu$; then the scale is $c_0/|y'-y|$, $c_0=2e^{-\gamma_E}$.)

\section{Line 2 does \emph{not} contribute to the running coupling}\label{sec:line2}
\paragraph{(a) Its $\mu$-dependence is zero for $\mathbf k\neq0$.} In $\LO^{(y)}_\xi$ the kernel $\frac{(x-y)\cdot(x'-y)}{(x-y)^2(x'-y)^2}$ and $U(y)$ do not depend on $y'$. The only $y'$ dependence is the phase and the weight. Write the weight as $\log(y'-y)^2+\log\frac{\mu^2}{4}$ (and $2\gamma_E$, $-\frac2\epsilon$ if added). The constant part gives
\begin{equation}
\Big[\log\frac{\mu^2}{4}+\text{const}\Big]\int d^2y'\,e^{-ik\cdot(y'-y)/\xi}=\Big[\log\frac{\mu^2}{4}+\text{const}\Big](2\pi)^2\xi^2\,\delta^2(\mathbf k)=0\qquad(\mathbf k\neq0).
\end{equation}
So line 2 has \textbf{no $\mu$}, \textbf{no pole}, and hence nothing to renormalize the coupling. Only $\log(y'-y)^2$ survives; from $\nabla^2\log r^2=4\pi\delta^2(r)$,
\begin{equation}
\int d^2y'\,e^{-ik\cdot(y'-y)/\xi}\log(y'-y)^2=-\frac{4\pi\xi^2}{\mathbf k^2}\qquad(\mathbf k\neq0),
\end{equation}
which is finite and $\mu$-independent.

\paragraph{(b) Its value.} With $\int_a^1d\xi\,\frac{1-\xi}{\xi^3}\,\xi^2=\log\frac{\vee}{k^+}-\big(1-\frac{k^+}{\vee}\big)$ and $\big|\int_w e^{ik\cdot w}\frac{(w-z)^i}{(w-z)^2}\big|^2=\frac{4\pi^2}{\mathbf k^2}$ (normalisations checked in \texttt{checks.py}),
\begin{equation}
\text{line 2}=-\Big[\log\frac{\vee}{k^+}-\Big(1-\frac{k^+}{\vee}\Big)\Big]H_4.\label{eq:l2}
\end{equation}
This is the $y'=y$ subtraction of the small-$\xi$ ($\frac{1-\xi}{\xi}$) part of line 1: line 1 contains $+\log\frac{\vee}{k^+}H_4$ from $\xi\to0$, and \eqref{eq:l2} removes it. That is a rapidity (BFKL) log, not a UV/coupling log. That is why it is kept inside the DGLAP integral in \eqref{eq:dglap}, where the two cancel and no $\log\frac{\vee}{k^+}$ is shown.

\paragraph{(c) Wrong structure for running coupling.} A running-coupling term must be $\propto\beta_0\,\delta(1-\xi)\times\LO\times\log\mu_R^2$: it multiplies LO itself, at $\xi=1$. Line 2 has the colour weight $2N_c\frac{1-\xi}{\xi}$, integrated over $\xi$, and multiplies $\LO^{(y)}_\xi$, not LO.

\paragraph{(d) Consistency.} Line 1 alone already gives the full $\mu$-dependence: $-\frac{\alpha_s}{2\pi}P_{gg}\otimes\LO$ (cancels DGLAP of $D$) and $+\frac{\alpha_s\beta_0}{4\pi}\LO$ (cancels the running of $\alpha_s$ in LO). A $\mu$-term from line 2 would spoil both.

\paragraph{Caveat.} Only if $\mathbf k$ is integrated over (so $\delta^2(\mathbf k)$ does not vanish) would line 2's constant part matter. For the measured $\mathbf k\neq0$ it does not.

\end{document}
